\chapter{تحلیل نیازمندی‌ها}

این فصل نیازمندی‌های پروژه را از زاویه تجربه کاربر و رفتار سامانه دسته‌بندی می‌کند. هدف این نیست که فقط یک فهرست خشک ارائه شود؛ هر نیازمندی باید نشان دهد کاربر چه انتظاری دارد و سامانه برای پاسخ به آن چه رفتاری باید داشته باشد.

\section{نقش‌های کاربری}
سه نقش اصلی در پروژه دیده می‌شود: مهمان، توسعه‌دهنده احرازشده و مدیر سیستم. مهمان بخش عمومی کاتالوگ را می‌بیند. توسعه‌دهنده بعد از ورود، با \lr{Dashboard}، انتشار \lr{API}، \lr{Caller} و \lr{Studio} کار می‌کند. مدیر سیستم هم مسئول کنترل کلی و نگهداری داده‌هاست.

\begin{figure}[htbp]
  \centering
  \begin{tikzpicture}[
    >=Latex,
    every node/.style={font=\fontsize{12}{14}\selectfont, align=center},
    actor/.style={draw=damghanNavy, very thick, rounded corners=2pt, fill=tablehead, minimum width=2.9cm, minimum height=0.75cm},
    action/.style={draw=tableline, thick, rounded corners=2pt, fill=white, minimum width=3.1cm, minimum height=0.72cm},
    flow/.style={-{Latex[length=2.2mm]}, thick, color=damghanTeal}
  ]
    \node[actor] (guest) at (4,2.2) {مهمان};
    \node[actor] (developer) at (4,0) {توسعه‌دهنده};
    \node[actor] (admin) at (4,-2.2) {مدیر سیستم};

    \node[action] (browse) at (-2.2,2.8) {مرور کاتالوگ};
    \node[action] (docs) at (-2.2,1.6) {مشاهده مستندات};
    \node[action] (account) at (-2.2,0.6) {حساب و سازمان};
    \node[action] (release) at (-2.2,-0.6) {انتشار \lr{API}};
    \node[action] (studio) at (-2.2,-1.8) {\lr{Caller} و \lr{Studio}};
    \node[action] (panel) at (-2.2,-3.0) {پنل مدیریت};

    \draw[flow] (guest.west) -- (browse.east);
    \draw[flow] (guest.west) -- (docs.east);
    \draw[flow] (developer.west) -- (account.east);
    \draw[flow] (developer.west) -- (release.east);
    \draw[flow] (developer.west) -- (studio.east);
    \draw[flow] (admin.west) -- (panel.east);
  \end{tikzpicture}
  \caption{نمای نقش‌ها و جریان‌های اصلی پروژه}
  \label{fig:roles}
\end{figure}

نقش مهمان بدون ورود می‌تواند بخش عمومی بازارچه را ببیند؛ این بخش برای جذب کاربر و معرفی \lr{API}ها ضروری است. نقش توسعه‌دهنده پس از احراز هویت به امکاناتی دسترسی دارد که با داده خصوصی کاربر ارتباط دارند، مانند کلید \lr{API}، سازمان، اشتراک، مصرف، انتشار و ابزارهای آزمایشی. نقش مدیر سیستم بیشتر به نگهداری داده و پشتیبانی از مسیرهای مدیریتی مربوط می‌شود و از نظر سطح دسترسی از دو نقش دیگر حساس‌تر است.

\section{نیازمندی‌های عملکردی و غیرفعال}
جدول \ref{tab:reqs} نیازمندی‌ها را در دو گروه اصلی نشان می‌دهد. نیازمندی‌های عملکردی به کارهایی مربوط‌اند که کاربر مستقیماً می‌بیند؛ مثل جستجوی \lr{API}، ورود، انتشار یا مشاهده مصرف. نیازمندی‌های امنیتی و اجرایی کیفیت همین رفتارها را حفظ می‌کنند؛ مثل پنهان‌کردن کلیدها، پاسخ خطای قابل فهم و اجرای محلی با \lr{Docker Compose}.

\begin{longtable}{p{1.35cm} p{5.25cm} p{1.75cm} p{1.35cm} p{1.85cm}}
\caption{نیازمندی‌های اصلی پروژه}\label{tab:reqs}\\
\toprule
\rowcolor{tablehead}\textbf{شناسه} & \textbf{نیازمندی} & \textbf{نوع} & \textbf{اولویت} & \textbf{وضعیت} \\
\midrule
\endfirsthead
\toprule
\rowcolor{tablehead}\textbf{شناسه} & \textbf{نیازمندی} & \textbf{نوع} & \textbf{اولویت} & \textbf{وضعیت} \\
\midrule
\endhead
\ModernTableStart
\lr{R1} & نمایش فهرست \lr{API}ها با صفحه‌بندی، جستجو، دسته‌بندی و مرتب‌سازی & عملکردی & زیاد & انجام‌شده \\
\lr{R2} & نمایش جزئیات، پلن‌ها، مستندات، \lr{endpoint}ها و \lr{API}های مشابه & عملکردی & زیاد & انجام‌شده \\
\lr{R3} & ثبت‌نام، ورود، نشست، خروج و کشف ارائه‌دهندگان ورود اجتماعی & عملکردی & زیاد & انجام‌شده \\
\lr{R4} & مدیریت کاربر، پروفایل، سازمان، اشتراک، پرداخت دستی و مصرف & عملکردی & زیاد & انجام‌شده \\
\lr{R5} & انتشار \lr{API} توسط کاربر احرازشده & عملکردی & متوسط & انجام‌شده \\
\lr{R6} & اجرای درخواست نمونه در \lr{Caller} و ثبت مصرف & عملکردی & متوسط & انجام‌شده \\
\lr{R7} & استقرار جریان \lr{Studio} و ثبت مصرف & عملکردی & متوسط & انجام‌شده \\
\lr{R8} & بازگشت خطای یکنواخت برای اعتبارسنجی و دسترسی & غیرفعال & زیاد & انجام‌شده \\
\lr{R9} & پنهان‌سازی کلیدها و توکن‌ها در پاسخ‌ها و لاگ‌ها & امنیتی & زیاد & انجام‌شده \\
\lr{R10} & اجرای محلی دو سرویس با \lr{Docker Compose} & استقرار & متوسط & مستند \\
\ModernTableStop
\bottomrule
\end{longtable}

اولویت زیاد برای نیازهایی در نظر گرفته شده که نبود آن‌ها باعث از کار افتادن هسته محصول می‌شود؛ مانند کاتالوگ، احراز هویت، حساب کاربری و امنیت داده حساس. اولویت متوسط برای قابلیت‌هایی است که ارزش محصول را کامل‌تر می‌کنند، اما می‌توانند در نسخه‌های بعدی نیز تکامل پیدا کنند؛ مانند \lr{Studio}، \lr{Caller} و اجرای محلی.

\section{ورودی‌ها، خروجی‌ها و اعتبارسنجی}
ورودی‌های اصلی شامل مشخصات حساب، پروفایل، شناسه یا نامک \lr{API}، فیلترهای جستجو، امتیاز، داده انتشار \lr{API}، شناسه پلن اشتراک و گره‌های \lr{Studio} است. در معماری مبتنی بر \lr{Django REST Framework}، معمولاً \lr{Serializer} مسئول کنترل و تبدیل همین داده‌هاست \cite{drf}.

خروجی‌های اصلی سامانه در قالب \lr{JSON} بازگردانده می‌شوند و برای مصرف فرانت‌اند طراحی شده‌اند. پاسخ‌های عمومی معمولاً شامل داده کاتالوگ، مستندات، پلن‌ها و وضعیت سلامت هستند. پاسخ‌های خصوصی شامل داده‌هایی مانند پروفایل، سازمان‌ها، کلید پنهان‌شده، اشتراک و تاریخچه مصرف هستند. جداسازی این دو دسته خروجی اهمیت دارد، زیرا داده عمومی برای همه کاربران قابل نمایش است، اما داده خصوصی باید فقط بعد از احراز هویت برگردد.

اعتبارسنجی در این پروژه دو نقش دارد: نخست جلوگیری از ذخیره داده ناقص یا نامعتبر، و دوم تبدیل خطا به پاسخ قابل فهم برای \lr{Client}. برای نمونه، هنگام ثبت امتیاز باید مقدار امتیاز معتبر باشد، هنگام انتشار \lr{API} فیلدهای اصلی باید حاضر باشند و هنگام اجرای \lr{Studio} ساختار گره‌ها باید قابل پردازش باشد.

\section{امنیت}
احراز هویت باید به شکلی باشد که کاربر بعد از ورود بتواند به بخش‌های خصوصی دسترسی داشته باشد، اما داده حساس او در پاسخ‌ها آشکار نشود. در پروژه‌هایی که با \lr{API Key}، توکن و کوکی کار می‌کنند، این موضوع فقط یک انتخاب ظاهری نیست؛ بخشی از امنیت اصلی محصول است \cite{django,drf}.

از دید نیازمندی، امنیت فقط ورود کاربر نیست؛ بلکه کنترل دسترسی، نگهداری نشست، جلوگیری از افشای رازها و محدودکردن خروجی‌های حساس را نیز شامل می‌شود. آزمون‌های موجود نشان می‌دهند که پنهان‌سازی داده حساس در خروجی‌ها بررسی شده است. با این حال، برای محیط تولید باید رازهای پیکربندی از فایل تنظیمات خارج شوند و سیاست‌های سخت‌گیرانه‌تری برای کوکی، دامنه مجاز و خطاهای عملیاتی اعمال شود.

\section{جمع‌بندی فصل}
نیازهای پروژه حول کشف عمومی \lr{API}، فضای کاری توسعه‌دهنده، مدیریت حساب و اشتراک، و امنیت نشست شکل گرفته‌اند. این نیازمندی‌ها مبنای فصل معماری هستند؛ چون لایه‌های سامانه باید همین نقش‌ها، جریان‌ها و محدودیت‌های امنیتی را پشتیبانی کنند.
